\section{Limitations}

As reflected in the quantitative results, CAC is by no means a perfect method for localized text-to-image generation. In this section, we discuss the limitations of our method and provide qualitative illustrations in Figure~\ref{fig:failure_cases}

The performance of our method heavily relies on the base model of choice. For example, since Stable Diffusion is only trained with English captioned images, CAC with Stable Diffusion also has limited performance on non-English prompts. Furthermore, when the base model fails to generate objects or features mentioned in the caption (as oppose to localized prompts), it is very difficult for our method to remedy that mistake. Moreover, similar to~\citet{sdedit}, although usually lying in a certain small range, the optimal hyperparameters such as the MD ratios and $\lambda$ are different for individual inputs. Sub-optimal hyperparameters can compromise the quality of the generated images. In practice, we encourage users to search for the best hyperparameters for their inputs. Our method also shows weaker performance with more complicated location information where the generated images tend to look less coherent and accurate to the location information.

\begin{figure}[t]
  \centering
  \includegraphics[width=\textwidth]{img/failure_cases.pdf}
  \caption{Example failure cases of our method. In general, the performance of our method depends of the base model and the selection of hyperparameters.}
  \label{fig:failure_cases}
\end{figure}